É possível construir painéis com os dados coletados em sala de aula, permitindo uma melhor visualização e consolidação das informações. Para essa finalidade, utilizou-se o Power BI para construir um \textit{dashboard} demonstrativo, ilustrado na Figura~\ref{fig:dashboard} abaixo.

\begin{figure}[H]
    \centering
    \begin{minipage}{0.999\textwidth}
        \centering
        \includegraphics[width=.9\textwidth]{espeon/espeon_dashboard.png}
        \caption{Power BI ESPEON}
        \label{fig:dashboard}
    \end{minipage}\hfill
\end{figure}

A ingestão de dados é realizada no editor Power Query M, apontando diretamente para a API produtiva, conforme apresentado na Figura~\ref{fig:espeon_query}.

\begin{figure}[H]
    \centering
    \begin{minipage}{0.999\textwidth}
        \centering
        \includegraphics[width=.9\textwidth]{espeon/bi_espeon_query.png}
        \caption{Consulta à API ESPEON}
        \label{fig:espeon_query}
    \end{minipage}\hfill
\end{figure}

Os dados são retornados como \textit{records}, os quais devem ser expandidos para que se convertam em tabelas. A expansão consiste em mapear as colunas originais --- isto é, as propriedades retornadas pela API --- a partir de \textit{Column1}, obtendo-se o conteúdo estruturado como tabela, conforme indicado nas Figuras~\ref{fig:espeon_query_load} e~\ref{fig:espeon_query_expand}.

\begin{figure}[H]
    \centering
    \begin{minipage}{0.999\textwidth}
        \centering
        \includegraphics[width=.9\textwidth]{espeon/bi_espeon_load.png}
        \caption{Carregamento de Dados}
        \label{fig:espeon_query_load}
    \end{minipage}\hfill
\end{figure}

\begin{figure}[H]
    \centering
    \begin{minipage}{0.999\textwidth}
        \centering
        \includegraphics[width=.9\textwidth]{espeon/bi_espeon_expand.png}
        \caption{Coluna Expandida}
        \label{fig:espeon_query_expand}
    \end{minipage}\hfill
\end{figure}
